% ════════════
% WP5_Working_Paper.tex
% Relational Dignity Infrastructure: Engineering the Street-to-Home Transition Protocol
% Material Dignity Infrastructure Working Paper 5
% ════════════

\documentclass[12pt, letterpaper]{article}

% ── PACKAGES ──
\usepackage[left=0.7in, right=0.7in, top=1in, bottom=1in, headheight=14pt]{geometry}
\usepackage[expansion=false]{microtype}
\usepackage{textcomp}
\usepackage{setspace}
\usepackage{booktabs}
\usepackage{tabularx}
\usepackage[titles]{tocloft}
\usepackage{tabularray}
\usepackage{longtable}
\usepackage{array}
\usepackage{multirow}
\usepackage{hyperref}
% natbib not used — manual thebibliography
\usepackage{amsmath}
\usepackage{amssymb}
\usepackage{parskip}
\usepackage{titlesec}
\usepackage{fancyhdr}
\usepackage[table]{xcolor}
\usepackage{caption}
\usepackage{footnote}
\usepackage{enumitem}
\usepackage{tikz}
\usepackage{needspace}
\usepackage{etoolbox}
\usepackage{float}

% ── COLOR SYSTEM ───
\definecolor{auditblue}{HTML}{003366}
\definecolor{rowgray}{HTML}{F5F5F5}

% ── GLOBAL SPACING ───
\setstretch{1.4}
\setlength{\parindent}{0pt}
\setlength{\parskip}{8pt}
\hyphenpenalty=1000

% ── SECTION FORMATTING. ZERO BOLDING IN TEXT. ───
\titleformat{\section}
  {\normalfont\large\color{auditblue}}{\thesection.}{0.5em}{}
\titleformat{\subsection}
  {\normalfont\normalsize\color{auditblue}}{\thesubsection.}{0.5em}{}
\titlespacing*{\section}{0pt}{18pt}{6pt}
\titlespacing*{\subsection}{0pt}{12pt}{4pt}

% ── HEADER / FOOTER ───
\pagestyle{fancy}
\fancyhf{}
\rhead{\small\textit{Relational Engineering Specifications (RDI Working Paper 5)}}
\lhead{\small\textit{Working Paper, June 2026}}
\cfoot{\thepage}
\renewcommand{\headrulewidth}{0.4pt}

% ── TABLE SETTINGS ───
\renewcommand{\arraystretch}{1.3}

% ── ORPHAN / WIDOW CONTROL ───
\clubpenalty=10000
\widowpenalty=10000
\displaywidowpenalty=10000

% ── BIBLIOGRAPHY ORPHAN PREVENTION ───
\AtBeginEnvironment{thebibliography}{\interlinepenalty=10000}

% ── HYPERREF CONFIGURATION ───
\hypersetup{
  colorlinks=true,
  linkcolor=black,
  citecolor=black,
  urlcolor=black,
  pdftitle={Relational Dignity Infrastructure: Engineering the Street-to-Home Transition Protocol},
  pdfauthor={DiBella, C.J.},
  pdfsubject={Housing First. Ontological Security. Identity Capital. Homelessness Recovery.},
  pdfkeywords={Relational Dignity Infrastructure, ontological security, identity capital, Housing First retention, Dunbar scale, Pod Steward, peer support, social network preservation, homelessness recovery, material dignity infrastructure, warm offer, institutional objectification, street-to-home pipeline}
}
\setcounter{tocdepth}{2}

\begin{document}
\captionsetup[table]{labelfont=normalfont, labelsep=period, skip=6pt}

% ── FRONT MATTER ───
\begin{titlepage}
\begin{center}
\setstretch{1.4}
{\LARGE \textbf{Relational Dignity Infrastructure}}\\[6pt]
{\large Engineering the Street-to-Home Transition Protocol}\\[0.5cm]
{\normalsize \textbf{Charles J. DiBella}}\\[0.01cm]
{\normalsize Principal Systems Architect}\\[0.01cm]
{\normalsize Working Paper --- June 2026}\\[0.3cm]
\textbf{ABSTRACT}
\vspace{0.2cm}
\end{center}
\begin{minipage}{0.95\textwidth}
\setstretch{1.15}
\noindent

Working Paper 4 of this series defined Relational Dignity Infrastructure (RDI) as the systematically designed social environment that produces ontological security, sustains identity capital accumulation, and maintains individual recognition within the MDI tower architecture. That paper identified two integration gaps requiring specification: the Relational Intake Protocol, grounding the warm offer in Porges's polyvagal neuroception architecture, and the Social Network Transition Protocol, operationalizing Social Network Analysis methodology for pod assignment. This paper executes both specifications and appends a third contribution: a systematic political economy of opposition analysis cataloging the eight objection categories the MDI-RDI model will encounter at deployment, with sourced counterpoints derived from the empirical and theoretical literature. The Relational Intake Protocol is specified as a three-state Autonomic Routing Matrix that determines approach vector, physical posture, and verbal script based on the observable autonomic state of the individual at the moment of contact, grounding the clinical principle that a refusal recorded in dorsal-vagal shutdown is a defense reflex, not a cognitive refusal of housing. The Social Network Transition Protocol is specified as a four-domain interview instrument, a reciprocity-weighted scoring algorithm, and a Process 4 pod assignment logic with four priority-ordered constraints. Together these specifications operationalize the two most consequential relational decisions in the street-to-home pipeline: who makes the offer and how, and who lives with whom after the offer is accepted. The political economy analysis establishes that the eight principal objection categories --- fiscal efficiency, market libertarian, NIMBY, Housing First fidelity, civil libertarian, incumbent provider, regulatory-permitting, and academic-methodological --- are each answerable from within the MDI-RDI evidence base, existing California law, and documented international operational outcomes, and that the cumulative architecture of the MDI model has been structured to reduce the specific friction that each objection category represents.

\end{minipage}
\end{titlepage}

\pagenumbering{arabic}

% ── DOCUMENT METADATA ────
\newpage
\section*{Document Metadata}

{\footnotesize\setstretch{1.25}

\noindent\textbf{Keywords}\\[1pt]
Relational intake protocol, polyvagal theory, autonomic routing, neuroception, social network analysis, pod assignment algorithm, warm offer, street-to-home pipeline, political economy of opposition, incumbent provider resistance, Housing First fidelity, compelled care, MDI, relational dignity infrastructure, Dunbar pod

\vspace{8pt}
\noindent\textbf{JEL Classification}
\begin{itemize}[nosep, before={\vspace{-8pt}}, leftmargin=2em]
    \item I14. Health and Inequality
    \item I38. Government Policy; Provision and Effects of Welfare Programs
    \item H75. State and Local Government: Health, Education, Welfare
    \item Z13. Economic Sociology; Economic Anthropology; Social and Economic Stratification
    \item R21. Urban, Rural, Regional, Real Estate, and Transportation Economics: Housing Demand
    \item D72. Political Processes: Rent-Seeking, Lobbying, Elections, Legislatures, and Voting
\end{itemize}

\vspace{8pt}
\noindent\textbf{Target eJournals}
\begin{itemize}[nosep, before={\vspace{-8pt}}, leftmargin=2em]
    \item Housing and Residential Economics eJournal
    \item Public Health Law and Policy eJournal
    \item State and Local Government eJournal
    \item Social Capital, Networks \& Trust eJournal
    \item Poverty, Income Distribution \& Social Protection eJournal
    \item Behavioral \& Experimental Economics eJournal
\end{itemize}

\vspace{8pt}
\noindent\textbf{Statement of Necessity}\\[1pt]
A theoretical framework without operational specifications is a research contribution without deployment capacity. Working Paper 4 named the relational layer and established its production conditions. This paper writes the two field manuals that the relational layer requires to move from theoretical specification to operational practice. The Autonomic Routing Matrix is the manual for the outreach worker on the street. The Social Network Analysis pipeline is the manual for the Process 4 data administrator managing pod assignments. Neither document existed before this paper. Both are required before the warm offer can be executed with relational fidelity at scale.

The political economy appendix exists because a deployment-oriented paper cannot present operational specifications without acknowledging the organized opposition those specifications will face. The academic literature on public goods provision (Olson, 1965), bureaucratic rent-seeking (Buchanan and Tullock, 1962), and high-modernist institutional failure (Scott, 1998) establishes that well-designed proposals consistently fail against concentrated incumbent resistance unless the proposal's deployment architecture explicitly addresses each resistance mechanism. This paper addresses each.

The field evidence grounding the relational intake specification is drawn in part from DiBella's own longitudinal street outreach documentation (500\_relational\_ethnography corpus), in which a fixed-schedule, needs-responsive, unconditional-regard outreach methodology produced documented trust accumulation over repeated encounters with chronically street-present individuals on Skid Row, Los Angeles. The autonomic state classification in the Routing Matrix is not a theoretical import from clinical psychology. It is a formalization of what that field practice learned empirically: that approach vector, physical posture, and verbal script must match the observable neurological state of the individual, not the institutional preference of the worker.

\newpage

\vspace{8pt}
\noindent\textbf{Author's Note}\\[1pt]
This paper is the fifth working paper in the Material Dignity Infrastructure series. MDI-D established that the relational layer must be designed to specification, built into the operational architecture, and measured against defined thresholds. This paper writes the specifications. MDI-D without MDI-E is a blueprint without a build manual. This paper provides the build manual for the two operational decisions that determine whether the relational layer functions as designed from day one of tower operation.

\vspace{8pt}
\noindent\textbf{Suggested Citation}\\[1pt]
DiBella, C.J. (2026). Relational Dignity Infrastructure: Engineering the Street-to-Home Transition Protocol. Material Dignity Infrastructure Working Paper 5. Working Paper, June 2026.

}

% ── TABLE OF CONTENTS ────
\newpage
\tableofcontents
\newpage
\setstretch{1.4}

\section{The Two Gaps This Paper Closes}

Working Paper 4 of this series defined Relational Dignity Infrastructure and identified two integration gaps between theoretical specification and operational practice. Gap 1: the warm offer lacks a relational specification --- it names the material elements (a room, a cart, a pet, a keycard) but not who makes the offer, in what manner, or under what autonomic conditions it is valid. Gap 2: the Social Network Transition Protocol, which the SNA literature establishes as the primary determinant of social network density at ninety days and therefore of twenty-four-month retention, does not yet exist as a field instrument or a pod assignment algorithm.

Both gaps sit at the same operational moment: the first hours of contact between the MDI system and a specific human being. The warm offer determines whether the street-to-home transition activates social engagement or threat defense in the individual's autonomic nervous system. The pod assignment algorithm determines whether the relational bonds that individual carries from street life survive the transition or are severed at the housing threshold. Both decisions are made once and produce consequences that accumulate across the entire residential tenure. Neither can be corrected after the fact without destroying the relational trust the system requires.

This paper closes both gaps and appends a third contribution: a systematic analysis of the eight objection categories the MDI-RDI architecture will encounter at deployment, with counterpoints grounded in existing California law, peer-reviewed evidence, and international operational benchmarks.

---

\section{The Field Evidence Base}

The specifications in this paper are not derived from theoretical inference alone. Longitudinal street outreach practice on Skid Row, Los Angeles --- documented in the relational ethnography corpus and summarized in the framework paper \textit{Relational Dignity as Infrastructure} (DiBella, 2026b) --- produced empirical observations that the Autonomic Routing Matrix formalizes.

The field methodology operated on a fixed weekly schedule at consistent locations. Material offerings --- tobacco, beverages, small currency --- were provided not as incentives but as signals of unconditional regard. The operative question in every encounter was: \textit{do you need anything?} Consistency of return was the primary trust-building mechanism. Where a request could not be fulfilled at the moment of contact, explicit commitment to future fulfillment was made and honored. This practice produced observable affective responses across multiple individuals --- characterized as documented written gratitude --- indicating that the experience of being treated as a full person, rather than as a service recipient, produced qualitatively distinct behavioral responses.

The critical operational finding: approach vector, physical posture, and verbal script cannot be fixed across all contact attempts. They must match the observable neurological state of the individual at the moment of contact. An individual in dorsal-vagal shutdown --- flat affect, immobility, dissociated gaze --- does not process questions. Asking "do you want to come inside?" at that moment is a category error. The question requires prefrontal processing capacity that the shutdown state has suspended. The offer fails not because the individual rejects housing but because the offer was made to a nervous system that cannot receive it.

Polyvagal theory (Porges, 2011) provides the neurobiological framework that explains this field observation. The Autonomic Routing Matrix translates that framework into an operational protocol.

---

\section{Gap 1: The Relational Intake Protocol}

\subsection{The Neuroception Architecture}

Porges's polyvagal theory identifies neuroception as the nervous system's non-conscious process of continuously scanning the environment for cues of safety or threat, operating below conscious awareness and prior to any behavioral response. Chronic street exposure locks neuroception into one of two defensive configurations. Dorsal-vagal shutdown --- the phylogenetically oldest defense state --- is characterized by immobility, metabolic conservation, flat affect, and dissociation. Sympathetic mobilization --- the fight-or-flight state --- is characterized by hypervigilance, rapid movement, elevated arousal, and rejection of proximity.

A warm offer delivered to either defensive state is structurally invalid. Refusal in dorsal-vagal shutdown is not a cognitive decision. It is an autonomic defense response that the prefrontal cortex did not generate and therefore cannot be attributed to the individual's considered preferences. Refusal in sympathetic mobilization is an adrenaline-driven rejection of proximity that will resolve when the physiological cycle completes. Recording either as a housing refusal and closing the application is an institutional category error that converts a neurological state into a policy outcome.

The Relational Intake Protocol begins with state identification. The worker evaluates observable biological indicators --- not behavior, not compliance, not presentation --- and routes the contact to the protocol appropriate to the identified state.

\subsection{The Autonomic Routing Matrix}

\textbf{State 1: Dorsal Vagal (Shutdown)}

\textit{Observable indicators:} Total or near-total immobility. Slumped or collapsed posture. Shallow or imperceptible breathing. Flat or absent affect. Vacant gaze. Apparent dissociation from immediate environment. Does not track the approach of the worker.

\textit{Threat perception:} Extreme or inescapable. The organism conserves metabolic energy by suspending voluntary activity.

\textit{Approach vector:} Oblique approach at minimum six feet. Do not approach head-on. Forbidden: direct sustained eye contact, which the shutdown nervous system reads as predation.

\textit{Physical posture:} Lower center of gravity to match the individual's level. Sit or crouch if the individual is on the ground. Hands visible and still.

\textit{Script constraint:} Forbidden: interrogative questions of any form. "How are you?" "Do you want to come inside?" "Can I help you?" All interrogatives demand prefrontal processing capacity that does not exist in State 1.

\textit{Mandatory script:} Declarative statements of safety and non-demand only. "I am going to leave this water here." "I am going to sit here for two minutes, and then I will leave." "I will come back on Thursday."

\textit{Operational routing:} Deliver material supply without requiring acknowledgment. Log the contact, the state, and the material provided. Withdraw. Repeat on the established schedule until autonomic shift to State 2 or State 3 is observed. Do not record as a service refusal.

\textbf{State 2: Sympathetic (Mobilized)}

\textit{Observable indicators:} Pacing or repetitive movement. Rigid musculature. Rapid chest-level breathing. Hyper-vigilant scanning of environment. Accelerated or pressured speech. Sweating.

\textit{Threat perception:} Active and immediate. The organism is flooded with adrenaline and prepared to fight or flee.

\textit{Approach vector:} Maintain distance. Never block an exit path. Position parallel to the individual, facing the same direction rather than face-to-face. This posture signals no predation intent.

\textit{Physical posture:} Standing, relaxed, hands open and visible. Slow all physical movements to approximately half normal speed.

\textit{Script constraint:} Forbidden: commands to calm down, logical arguments about the offer, or any statement that requires the individual to evaluate a choice. Adrenaline forecloses deliberation.

\textit{Mandatory script:} Match the energy level but lower the pitch. Short concrete sentences. Validate the threat perception without confirming delusion. "You are moving. I am staying right here. You have space to move."

\textit{Operational routing:} Allow the adrenaline cycle to complete --- typically five to fifteen minutes of sustained non-escalation. Do not introduce the housing offer until physiological deceleration is observable: breathing slows, pacing stops, scanning frequency decreases. Record the contact and the state.

\textbf{State 3: Ventral Vagal (Engaged)}

\textit{Observable indicators:} Relaxed or open posture. Deep belly breathing. Spontaneous eye contact at normal duration. Tonal variation in voice. Facial micro-expressions of recognition or interest.

\textit{Threat perception:} Absent or low. The organism is capable of social engagement and prefrontal cognitive processing.

\textit{Approach vector:} Direct approach. Face-to-face orientation is appropriate. Normal proximity.

\textit{Physical posture:} Relaxed, mirroring the individual's posture. Normal eye contact.

\textit{Script constraint:} The prefrontal cortex is online. Choices, logistics, and future planning are biologically accessible.

\textit{Mandatory script:} Execute the full warm offer. Reference specific historical knowledge to demonstrate relational continuity: "I know you need Ranger to come with you. We have a secure kennel on the third floor with his name on it. Your cart is already stored. The room has a lock and the key is yours." Specificity is the evidence of relationship.

\textit{Operational routing:} Proceed to pod transition logistics. The offer is valid.

\subsection{The Intake Failure Rule}

An intake offer made while the individual is in State 1 or State 2 is structurally invalid regardless of what response is recorded. Refusal in State 1 or State 2 is not a refusal of housing. It is an autonomic defense response. The offer must be withheld or withdrawn until State 3 is achieved through sustained co-regulation --- the process by which the worker's own ventral vagal state, maintained through calm, non-demanding presence, gradually shifts the individual's neuroception toward safety.

Recording a State 1 or State 2 response as a housing refusal and closing the application is an institutional falsification of the individual's actual decision-making capacity. The MDI operational protocol prohibits this recording and requires the contact to be logged as a state observation with a scheduled follow-up.

\subsection{The Non-Expiring Offer Protocol}

The warm offer does not expire. The individual who declines in State 3 receives the same offer at the next scheduled contact. The word "decline" is not used in MDI operational documentation because it implies a considered decision that the protocol cannot assume was made under conditions of full neurological availability. The offer is ongoing. The relationship is the delivery mechanism. The offer becomes credible over time precisely because it continues to arrive without consequence for previous responses.

---

\section{Gap 2: The Social Network Transition Protocol}

\subsection{The Retention Stakes}

Tsai and colleagues' 2012 longitudinal study across three US Housing First sites established social network density at ninety days as a statistically significant independent predictor of twelve-month housing stability, controlling for substance use, psychiatric diagnosis, and prior housing history. Each additional named reciprocal peer contact at ninety days added measurably to retention probability in a dose-response relationship. The social integration evidence is not correlational background. It is the primary mechanism by which the RDI layer produces the retention outcomes the MDI Singular Prototype Threshold measures.

Encampment communities contain existing social networks with documented reciprocal relational bonds. Housing assignment processes that ignore these networks --- through lottery, queue position, or geographic availability --- sever them at the housing threshold and convert a relational asset into a retention liability. Pod assignment is a clinical decision. The MDI system treats it as one.

\subsection{The Assessment Instrument}

The Social Network Analysis (SNA) assessment is conducted by the outreach team during the Phase Zero encampment engagement period, prior to tower activation. It uses material reliance questions rather than subjective friendship questions, because self-reported friendship is unreliable under conditions of chronic stress and survival competition. Material reliance --- who watches your gear, who you share food with, who you find first in a crisis --- is a behavioral record that can be independently verified.

\textbf{Mandatory Interview Script}

\begin{table}[h!]
\renewcommand{\arraystretch}{1.5}
\begin{tabularx}{\textwidth}{@{}lXl@{}}
\toprule
\textbf{Domain} \& \textbf{Question} \& \textbf{Target Output} \\
\midrule
\textbf{Material Security} \& "When you leave your tent, who watches your gear?" \& Identifies primary trust node \\
\textbf{Metabolic Supply} \& "Who do you share food or tobacco with?" \& Identifies routine reciprocal nodes \\
\textbf{Crisis Response} \& "If you get sick or hurt, who do you find first?" \& Identifies apex reliance node \\
\textbf{Toxicity Identification} \& "Is there anyone here you need to stay away from?" \& Identifies predator or parasite nodes requiring separation \\
\bottomrule
\end{tabularx}
\end{table}

All responses are logged by name into the Process 4 cybernetic grid. Unnamed responses are discarded.

\subsection{Verification and Scoring}

All claimed ties require bidirectional verification. The worker asks Node A if they watch Node B's gear. Then asks Node B if Node A watches their gear. Unverified claims are algorithmically weighted at zero and excluded from assignment logic. This prevents socially dominant individuals from manufacturing relational claims that inflate their pod placement leverage.

| Metric | Verification Action | Algorithmic Weight |
| :--- | :--- | :--- |
| \textbf{Reciprocity (Verified)} | Both nodes confirm the tie independently | \textbf{+3} (Primary Cluster Anchor) |
| \textbf{Reciprocity (Unverified)} | One node denies or is unavailable | \textbf{0} (Discard) |
| \textbf{Density (Hub)} | Node named by three or more individuals in the encampment | \textbf{+2} (Central Hub) |
| \textbf{Toxicity (Risk)} | Node named as threat by any verified node | \textbf{$-$5} (Isolation Flag) |

\subsection{Process 4 Assignment Algorithm}

The scored data enters the Process 4 cybernetic grid. Pod assignment executes the following four constraints in strict priority order:

\textbf{Priority One --- Isolation Constraint:} Any node carrying a Toxicity Flag ($-$5) is assigned to a physically separate floor or building from every node that named them as a threat. This constraint overrides all others. No exception.

\textbf{Priority Two --- Anchor Constraint:} All mutually verified Reciprocity pairs (+3) are mandatorily assigned to the same twelve-person Dunbar sub-pod within their floor. Verified bond partners who cannot be co-assigned within capacity constraints trigger a capacity expansion review before any separation is recorded.

\textbf{Priority Three --- Hub Constraint:} High-Density nodes (+2) are assigned to pods identified as requiring social stabilization --- pods with a high proportion of individuals carrying zero verified ties --- where their existing social gravity serves the pod's community formation rather than concentrating social capital in already-dense clusters.

\textbf{Priority Four --- Remnant Constraint:} Individuals with zero verified ties are distributed evenly across high-density pods. No pod may contain a majority of zero-tie individuals. Social isolation is a contagious environmental condition; concentrating isolated individuals produces a pod environment in which no community formation is possible regardless of physical design quality.

\subsection{The Assignment Failure Rule}

Assigning mutually verified bond partners (+3) to separate floors constitutes an architectural failure. The social network transition protocol identifies it as equivalent to offering a housing unit without a lock: the material infrastructure is present but the relational layer has been destroyed at intake. A pod assignment log that separates verified bond partners without a documented capacity justification and a documented escalation path for review is an operational violation of the RDI retention standard.

---

\section{Political Economy of Opposition: Compressed Counterpoint Framework}

The eight objection categories cataloged in full in the companion research document (DiBella, 2026c) are summarized here with their primary structural counterpoints. Each objection is stated at its maximum force. Each counterpoint is drawn from existing law, peer-reviewed evidence, or documented operational outcome.

\textbf{1. Fiscal and Efficiency Critics} argue that per-unit costs exceed alternatives and that CalAIM billing projections are undemonstrated. Counterpoint: the correct comparator is total annual street-level cost --- emergency department utilization, psychiatric holds, incarceration cycling, outdoor operations --- against which Culhane et al. (2002) document \$16,282 per person per year in avoided costs for Housing First participants and the HUD brief documents \$46,000 per year for heavy service users in permanent supportive housing. The CalAIM billing pathways (ECM, ILOS) operate under existing managed care contracts; they do not require new legislation.

\textbf{2. Market Libertarian Critics} argue that deregulation would solve housing scarcity without new institutional apparatus. Counterpoint: filtering theory (Rosenthal, 2014) establishes that the filtering mechanism does not reach a population whose neurological incapacity prevents market participation at any price. The MDI architecture uses market mechanisms --- commercial real estate conversion below replacement cost, performance-based retention incentives, Medi-Cal billing revenue --- to address a market-failure population. The techno-optimist critique and the MDI proposal address different problems within the same crisis.

\textbf{3. NIMBY Opposition} argues that siting 2,000 high-acuity individuals in a commercial corridor imposes measurable costs on adjacent stakeholders. Counterpoint: California AB 2011 (2022) creates by-right administrative review for qualifying commercial conversions, eliminating the discretionary review process through which neighborhood opposition typically mobilizes. The property value literature (Nguyen, 2005; Lee and Song, 2019) attributes negative effects to facility management quality, not facility type --- and the MDI management architecture is specifically designed to produce the management quality that the literature identifies as the variable. The ground floor commons civic permeability model converts the facility from threat to neighborhood asset.

\textbf{4. Housing First Purists} argue that congregate models violate Housing First fidelity standards. Counterpoint: the Finnish Y-Foundation model, which the purist literature cites as the global benchmark, uses congregate supported housing as a transitional layer within the Housing First pipeline. The MDI tower is Phase 1; the Tenancy Bridge Guarantee commits to permanent independent housing as Phase 2. The fidelity objection mistakes sequencing for violation.

\textbf{5. Civil Libertarian Critics}, following Szasz (1984), argue that compelled care through the legal pathway reproduces the asylum model under a new institutional form. Counterpoint: Szasz defeats unconstrained therapeutic state authority. He does not defeat a narrow, conjunctive, court-supervised, time-limited, reversible clinical intervention whose precondition is verified neurological destruction of the biological substrate of autonomous choice. The Anti-Detention Covenant (unrestricted egress for voluntary residents at all times) and the Rapid Transition Protocol (crisis produces transfer, not eviction) are the structural guarantees that the legal pathway remains exceptional rather than operational default.

\textbf{6. Incumbent Service Providers} will oppose through procurement gatekeeping, licensing delays, and lobbying for competing funding allocations. Olson (1965) predicts this precisely: concentrated interests with organizational capacity act against dispersed beneficiaries who cannot coordinate in proportion to their aggregate welfare stake. Counterpoint: the MDI funding architecture routes through private capital acquisition and Medi-Cal managed care billing rather than county homelessness services procurement contracts. The incumbent industry's political leverage depends on government funding dependence; the MDI model reduces that dependence structurally.

\textbf{7. Regulatory and Permitting Critics} argue that commercial conversion triggers ADA, seismic, fire suppression, and CEQA requirements that make timelines and costs unmanageable. Counterpoint: AB 2011 and SB 6 (2022) create by-right authorization for qualifying commercial conversion projects, converting 18-36 month discretionary review to 90-180 day administrative review. The STC 65 acoustic specification is defensible as an ADA reasonable accommodation for residents with PTSD-associated sensory processing conditions, converting a cost objection into a compliance obligation.

\textbf{8. Academic and Methodological Critics} argue that the MDI-RDI architecture lacks RCT-level validation. Counterpoint: the existing system --- 24 billion dollars of spending producing population growth --- has no RCT validation for its dominant service model. The MDI evidence synthesis (ontological security theory, polyvagal neuroscience, SNA methodology, identity capital research, Finnish operational outcomes) is stronger than the evidence base for most interventions the critique accepts unchallenged. The Singular Prototype Threshold creates a replication standard requiring demonstrated outcome at One California Plaza before network expansion --- the opposite of how government programs typically scale.

\textbf{Cross-cutting: The Deployment Problem.} Olson's collective action logic predicts that dispersed beneficiaries will not organize against concentrated incumbent resistance. Scott (1998) documents the consistent failure of high-modernist institutional schemes that ignore informal social knowledge. The MDI architecture's cumulative response: by-right permitting reduces regulatory gatekeeping, private capital acquisition reduces procurement dependence, CalAIM billing aligns managed care plan incentives, and the Singular Prototype Threshold prevents premature scale. No single feature guarantees deployment. Together they reduce each friction that has historically defeated comparable proposals.

---

\section{Integration with WP4 Verification Metrics}

The five RDI leading indicators established in WP4 acquire operational meaning from the specifications in this paper. Three of the five are directly traceable to the protocols above.

The \textbf{Named Peer Contact Index} ($\ge$3 reciprocal contacts per resident by day 90) is a direct output of the Social Network Transition Protocol. Pods assembled through the Anchor and Hub constraints arrive with pre-existing verified ties, accelerating the accumulation toward the three-contact threshold. A pod assembled through lottery arrives with zero verified ties and must build from nothing inside the tower environment.

The \textbf{Voluntary Service Engagement Rate} ($\ge$70\% by day 90) is a direct output of the Relational Intake Protocol. Residents who entered through a State 3 warm offer delivered by a known person with specific historical knowledge arrive with an existing institutional trust foundation. Residents who entered through a State 1 or State 2 offer that was improperly validated arrive with institutional trust at zero or below zero and must rebuild from a trust deficit.

The \textbf{Sustained Ground Floor Opt-Out Rate} ($\le$15\% over 90 days) is partially traceable to pod assignment quality. Residents with verified relational bonds assigned to their pod have social reasons to leave their unit and enter common spaces. Isolated residents assigned to pods where they know no one have no social incentive for common space engagement and are precisely the population the $\le$15\% threshold is designed to detect as a system signal.

The fiscal calculus from WP4 (Section VII) holds: the Named Peer Contact floor generates \$338,000 to \$455,000 in avoided processing costs per tower per year across thirteen pods, before any Medi-Cal billing is applied. The Social Network Transition Protocol is the operational mechanism that produces this metric. The warm offer relational specification is the operational mechanism that determines whether the resident arrives capable of achieving it.

---

\section{Contribution and Sequence}

This paper makes three contributions to the MDI-RDI series.

It converts theoretical production conditions into operational field protocols. WP4 established that the warm offer must be relational from its first moment and that social network bonds must be preserved across the housing transition. This paper specifies how both requirements are operationally satisfied at the level of approach vector, verbal script, scoring algorithm, and assignment constraint. Specification is what makes replication possible. Without it, each implementation team improvises, and improvisation produces the variable quality that defeats the Singular Prototype Threshold's replication logic.

It introduces field evidence from sustained street outreach practice as an empirical grounding for the Autonomic Routing Matrix. The matrix is not a theoretical import from clinical psychology applied speculatively to a homeless population. It is a formalization of what direct practice on Skid Row demonstrated across multiple individuals over an extended period: that the biological state of the individual at the moment of contact determines whether the contact produces trust accumulation or trust damage, regardless of the skill or goodwill of the worker.

It preempts the eight objection categories the model will face at deployment by answering each from within existing law, peer-reviewed evidence, and documented operational benchmark. Proposals that do not preempt their objections invite those objections to arrive as surprises. The MDI deployment architecture is designed to make surprises rare by resolving each known friction point before it becomes an obstacle.

---

\addcontentsline{toc}{section}{References}
\section*{References}

Buchanan, James M., and Gordon Tullock. \textit{The Calculus of Consent: Logical Foundations of Constitutional Democracy}. Ann Arbor: University of Michigan Press, 1962.

California Legislative Information. \textit{Assembly Bill 2011: Affordable Housing and High Road Jobs Act of 2022}. Sacramento, 2022.

California Legislative Information. \textit{Senate Bill 6: Middle Class Housing Act of 2022}. Sacramento, 2022.

Community Guide. "Housing First: Community Living for Adults Experiencing Homelessness." \textit{The Community Preventive Services Task Force}, 2020.

Culhane, Dennis P., Stephen Metraux, and Trevor Hadley. "Public Service Reductions Associated with Placement of Homeless Persons with Severe Mental Illness in Supportive Housing." \textit{Housing Policy Debate} 13, no. 1 (2002): 107–163.

DiBella, Charles J. \textit{Material Dignity Infrastructure: A Distributed Stewardship Model for Homelessness Resolution}. SSRN Working Paper 5968756. 2025.

DiBella, Charles J. \textit{Material Dignity Infrastructure: Structural Misalignment and the Activation of Surplus Shelter Capacity}. SSRN Working Paper 6211658. February 2026.

DiBella, Charles J. \textit{Material Dignity Infrastructure: Los Angeles Metropolitan Stabilization --- A Street-to-Home Pipeline Analysis}. SSRN Working Paper 6579600. May 2026.

DiBella, Charles J. \textit{Relational Dignity Infrastructure: The Human Layer Required to Make Material Housing Inhabitable}. SSRN Working Paper 6881539. June 2026.

DiBella, Charles J. "Relational Dignity as Infrastructure: A Field-Based Framework for Sustainable Homeless Outreach." Working paper. 2026b.

DiBella, Charles J. \textit{Political Economy of Opposition: Systematic Objections to the MDI-RDI Architecture and Counterpoint Framework}. MDI Research Document 10. June 2026c.

Grossman, Peter, and Edwin W. Taylor. "A Clarification and Defense of the Theory of Polyvagal Control." \textit{Biological Psychology} 74, no. 2 (2007): 260–268.

Lee, Sugie, and Yena Song. "The Effect of Supportive Housing on Surrounding Property Values." \textit{Housing Policy Debate} 29, no. 1 (2019): 162–176.

Nguyen, Mai Thi. "Does Affordable Housing Detrimentally Affect Property Values? A Review of the Literature." \textit{Journal of Planning Literature} 20, no. 1 (2005): 15–26.

Olson, Mancur. \textit{The Logic of Collective Action: Public Goods and the Theory of Groups}. Cambridge: Harvard University Press, 1965.

Porges, Stephen W. \textit{The Polyvagal Theory: Neurophysiological Foundations of Emotions, Attachment, Communication, and Self-Regulation}. New York: Norton, 2011.

Rosenthal, Stuart S. "Are Private Markets and Filtering a Viable Source of Low-Income Housing?" \textit{American Economic Review} 104, no. 2 (2014): 687–706.

Scott, James C. \textit{Seeing Like a State: How Certain Schemes to Improve the Human Condition Have Failed}. New Haven: Yale University Press, 1998.

Szasz, Thomas. \textit{The Therapeutic State: Psychiatry in the Mirror of Current Events}. Buffalo: Prometheus Books, 1984.

Tsai, Jack, Richard Stanhope, Robert A. Rosenheck, and Wesley J. Kasprow. "The Role of Housing, Substance Abuse, and Mental Illness in Community Participation and Community Integration." \textit{Psychiatric Services} 63, no. 5 (2012): 435–442.

Tsemberis, Sam. \textit{Housing First: The Pathways Model to End Homelessness for People with Mental Illness and Addiction}. Center City: Hazelden, 2010.

Wasserman, Stanley, and Katherine Faust. \textit{Social Network Analysis: Methods and Applications}. Cambridge: Cambridge University Press, 1994.

Y-Foundation. \textit{A Home of Your Own: Housing First and Ending Homelessness in Finland}. Helsinki: Y-Foundation, 2022.

\end{document}
